% ============================================================
%  Page de garde
% ============================================================
\begin{titlepage}
\begin{center}

% En-tête université
{\large\bfseries République Algérienne Démocratique et Populaire}\\[0.3em]
{\large Ministère de l'Enseignement Supérieur et de la Recherche Scientifique}\\[1.5em]

\rule{\linewidth}{0.6pt}\\[0.5em]
{\large\bfseries Université des Sciences et de la Technologie Houari Boumediene}\\
{\normalsize Faculté d'Informatique · Département SI}\\[0.3em]
\rule{\linewidth}{0.6pt}

\vspace{2em}

% Discipline
{\color{usthbblue}\large Master 1 Systèmes d'Information et d'Intelligence Artificielle}\\[0.5em]
{\normalsize Module : \textbf{RCR1 — Représentation des Connaissances et Raisonnement}}

\vspace{3em}

% Titre principal
\begin{tcolorbox}[defbox, width=0.92\linewidth, halign=center]
  {\LARGE\bfseries Implémentation de Formalismes de}\\[0.4em]
  {\LARGE\bfseries Représentation des Connaissances}\\[0.4em]
  {\LARGE\bfseries et Raisonnement}\\[1em]
  {\large Logique Modale · Logique des Défauts · Logique Classique}\\
  {\large Logique de Description · Réseaux Sémantiques}
\end{tcolorbox}

\vspace{3em}

% Domaines d'application
{\normalsize Domaines d'application :
  \textbf{Industrie musicale algérienne} et \textbf{Système judiciaire}}

\vspace{3em}

% Auteurs
\begin{tabular}{rl}
  \textbf{Réalisé par :} & D. Arslan \\
                          & [Prénom Nom Binôme 2] \\[1em]
  \textbf{Encadré par :} & [Nom de l'encadrant] \\[1em]
  \textbf{Groupe :}      & [Numéro du groupe]
\end{tabular}

\vfill

% Outils
{\small\color{gray}
  TweetyProject 1.27 · OWL API 4.5 · HermiT 1.4.3 · JavaFX 17 · Java 11 · Maven
}\\[1em]

\rule{\linewidth}{0.4pt}\\[0.4em]
{\normalsize Année universitaire 2025 -- 2026}

\end{center}
\end{titlepage}
